% cv_teaching_cjnowacek.tex
% Curriculum vitae for adjunct teaching applications (character rigging).
% Standalone document: no profile toggles, always builds the same PDF.
\documentclass[10pt,a4paper,sans]{moderncv}

\moderncvstyle{classic}
\moderncvcolor{green}

\usepackage[utf8]{inputenc}
\usepackage[scale=0.95]{geometry}
\usepackage{enumitem}
\setlength{\hintscolumnwidth}{2.8cm}
\setlength{\separatorcolumnwidth}{0.02\textwidth}
\setlength{\parskip}{0pt}

% ---------------------------
% PERSONAL INFO
% ---------------------------
\name{CJ}{Nowacek}
\title{Character Rigging \& Technical Art}
\address{Cleveland Heights, Ohio}{}{}
\phone[mobile]{(216) 469 0959}
\email{cj@cjnowacek.com}
\homepage{cjnowacek.com}
\social[github]{cjnowacek}

\begin{document}
\makecvtitle

% ---------------------------
% PROFILE
% ---------------------------
\section{Profile}
\cvitem{}{Adjunct faculty at the Cleveland Institute of Art, co-teaching 3D Rigging \& Problem Solving (ANIM 312) and advising BFA seniors on their senior projects and portfolios. Character rigger and technical artist who shipped \textbf{SMITE 2} at Hi-Rez Studios, with ongoing freelance rigging work and a full-time pipeline development role at an architectural visualization studio, where training production staff on in-house tools is a regular part of the job. Production practice spans rigging, skinning, and runtime optimization in Maya, 3ds Max, and Unreal Engine 5, alongside the Python tooling that supports artists.}

% ---------------------------
% EDUCATION
% ---------------------------
\section{Education}
\cventry{2023}{BFA, Animation and Game Design}{Cleveland Institute of Art}{Cleveland, OH}{}{Senior Project: \textit{Runaway}, a cinematic story-driven game}

% ---------------------------
% ACADEMIC APPOINTMENTS
% ---------------------------
\section{Academic Appointments}
\cventry{2026--Present}{Adjunct Faculty, Animation}{Cleveland Institute of Art}{Cleveland, OH}{}{
	\begin{itemize}[leftmargin=0.6cm, itemsep=0pt]
		\item Co-teach ANIM 312, 3D Rigging \& Problem Solving: a BFA studio course in Maya covering advanced character rigging (IK/FK switching, stretchy limbs, auto-twist and IK/FK spines, ribbon controls, corrective blendshapes, facial controls), dynamics and simulation, and motion-capture retargeting onto rigs, with graded exercises and presented projects
		\item Advise BFA seniors on senior projects, portfolios, and demo reels, and on preparing for technical-art and animation roles
	\end{itemize}}

% ---------------------------
% TEACHING, MENTORSHIP & KNOWLEDGE TRANSFER
% ---------------------------
\section{Teaching, Mentorship \& Knowledge Transfer}
\cvitem{Advising}{Advise BFA seniors at the Cleveland Institute of Art on senior projects, portfolios, and demo reels, and on the transition into technical-art and animation roles}
\cvitem{Staff Training}{Conducts training for production personnel at MediaLab 3D Solutions as a regular part of the pipeline role: initial walkthroughs when in-house tools are deployed and ongoing retraining as they are updated, with documentation authored for every tool}
\cvitem{Peer Instruction}{Helped fellow students build and debug character rigs in shared BFA courses at the Cleveland Institute of Art, providing informal one-on-one instruction}
\cvitem{Peer Review}{Participated in Hi-Rez Studios' regular peer-review process for character assets, giving and receiving structured critique on rigs, skinning, and tooling}
\cvitem{Technical Writing}{Publishes technical breakdowns of production rigging tools at cjnowacek.com, including a skin-weight transfer utility (Envelope Export Tool) and an automated FK chain-rigging tool built for \textbf{SMITE}}
\cvitem{Documentation}{Wrote documentation and runbooks for automated pipeline processes at Hi-Rez Studios, supporting team scaling and knowledge transfer across the technical animation team}
\cvitem{Open Source}{Maintains \textbf{maya-tools}, a public repository of Python rigging tools for Maya (github.com/cjnowacek/maya-tools)}
\cvitem{User Guides}{Authored end-user documentation for automated asset-processing tools at Anuttacon, incorporating stakeholder feedback to meet production needs}
% TODO (CJ): add real entries here as they happen, e.g.:
% \cvitem{Guest Lectures}{...classroom visits at other programs, portfolio reviews, demo sessions...}
% \cvitem{Service}{...committee work, faculty governance, curriculum development...}
% \cvitem{Workshops}{...rigging workshops, game jam instruction...}

% ---------------------------
% PROFESSIONAL EXPERIENCE
% ---------------------------
\section{Professional Experience}

\cventry{2026--Present}{CGI Pipeline Developer}{MediaLab 3D Solutions}{Remote}{}{
	\begin{itemize}[leftmargin=0.6cm, itemsep=0pt]
		\item Designed and built a Python scene-automation framework for scene cleanup, optimization, diagnostics, and render-farm submission, accessible via CLI, PyQt desktop app, and native 3ds Max integration
		\item Built artist-facing tools for the studio's central 3ds Max toolset, modernized its legacy MaxScript, and own the build-and-release workflow that ships it to production
		\item Develop test-first with AI coding agents: automated checks precede every change and a gated review precedes every merge
		\item Author the documentation for in-house tools and conduct training for production staff: initial walkthroughs at deployment and ongoing retraining as tools are updated
	\end{itemize}}

\cventry{2026}{MaxScript Specialist (Contract)}{MediaLab 3D Solutions}{Remote}{}{
	\begin{itemize}[leftmargin=0.6cm, itemsep=0pt]
		\item Built MaxScript and Python automation for asset processing and repetitive scene operations in 3ds Max; converted to full-time pipeline developer role
	\end{itemize}}

\cventry{2025}{Technical Artist (Contract)}{Anuttacon via GoDemic}{Remote}{}{
	\begin{itemize}[leftmargin=0.6cm, itemsep=0pt]
		\item Developed automated asset-processing tools in Python and MEL, iterating on feedback to improve usability and reliability
		\item Wrote user documentation and incorporated stakeholder feedback to meet production needs
	\end{itemize}}

\cventry{2024--2025}{Technical Animator (Contract)}{Hi-Rez Studios}{Remote}{}{
	\begin{itemize}[leftmargin=0.6cm, itemsep=0pt]
		\item Built Python and MaxScript automation for asset integration and deployment in an Unreal Engine 5 environment, eliminating manual setup and reducing deployment errors
		\item Collaborated with Animation, Design, and QA to find bottlenecks and automate milestone-delivery steps
	\end{itemize}}

\cventry{2023--Present}{Rigging Artist (Freelance)}{Independent}{Remote}{}{
	\begin{itemize}[leftmargin=0.6cm, itemsep=0pt]
		\item Character rigging for original characters and independent projects in Maya, including custom rigs for \textit{Grimmlins Tale}
	\end{itemize}}

\cventry{2023--2024}{Associate Technical Animator}{Hi-Rez Studios}{Alpharetta, GA}{}{
	\begin{itemize}[leftmargin=0.6cm, itemsep=0pt]
		\item Shipped \textbf{SMITE 2} as part of the technical animation team: character rigging alongside the Python/MaxScript tooling that supported the asset pipeline
		\item Built and refined character rigs, partnering with animation to deliver production-ready assets
		\item Optimized character assets for runtime performance in a live multiplayer title, keeping rigs and skinned meshes within performance budgets
		\item Developed automated batch-processing tools in Python and MaxScript, cutting manual asset-processing time by \textbf{25\%}
		\item Modernized legacy pipeline scripts that had gone years without maintenance, keeping the rigging team's daily tooling current and reliable
	\end{itemize}}

\cventry{2022--2023}{Technical Animator Intern}{Hi-Rez Studios}{Remote}{}{
	\begin{itemize}[leftmargin=0.6cm, itemsep=0pt]
		\item Developed Python scripts to automate repetitive production tasks, improving efficiency and reducing manual errors
	\end{itemize}}

% ---------------------------
% SHIPPED WORK & CREDITS
% ---------------------------
\section{Shipped Work \& Credits}
\cvitem{2024}{\textbf{SMITE 2} (Hi-Rez Studios): character rigging and asset-pipeline tooling for a live multiplayer title}
\cvitem{2023--Present}{\textit{Grimmlins Tale} and independent projects: freelance character rigs in Maya}

% ---------------------------
% EXHIBITIONS & SELECTED ACADEMIC PROJECTS
% ---------------------------
\section{Exhibitions \& Selected Academic Projects}
\cvitem{2023}{\textbf{Runaway}: story-driven game built solo in Unreal Engine 5 as BFA senior project, shown in the Cleveland Institute of Art senior exhibition; complete character rigging pipeline in Maya with facial animation and IK/FK blending, plus Python asset automation}
\cvitem{2022}{\textbf{The Sintern}: animated short film made with a four-person team; character rigging with custom facial and cartoon-style expression controls, modeling, lighting, and pipeline support in Maya}
% ---------------------------
% PROFESSIONAL DEVELOPMENT
% ---------------------------
\section{Professional Development}
\cvitem{Coursework}{The Gnomon Workshop: Character Rigging in Maya for Game Production; Rigging the Jaw with Python in Maya; Master Rigging \& Python Scripting in Maya}
\cvitem{Ongoing}{Self-directed study through production work, public tool development, and personal rigging projects}

% ---------------------------
% TECHNICAL SKILLS
% ---------------------------
\section{Technical Skills}
\cvitem{Rigging}{Character rigging and skinning in Maya and 3ds Max; deformation setup; rig and skinned-mesh optimization for real-time engines}
\cvitem{Programming}{Python (Maya API, PyMEL, PyQt/PySide, pymxs), MaxScript, MEL, Bash, Git, Perforce}
\cvitem{DCC/Engines}{Maya, 3ds Max, Unreal Engine 5, Unity, Substance Painter}

% ---------------------------
% SECTIONS TO ADD WHEN THERE IS REAL MATERIAL
% (standard academic-CV sections; uncomment and fill as they become true)
% ---------------------------
% \section{Talks \& Presentations}
% \cvitem{20XX}{...}
%
% \section{Professional Affiliations}
% \cvitem{}{...SIGGRAPH, IGDA, local game-dev groups...}
%
% \section{Awards \& Honors}
% \cvitem{20XX}{...}

% ---------------------------
% REFERENCES
% ---------------------------
\section{References}
\cvitem{}{Professional references from Hi-Rez Studios and MediaLab 3D Solutions available upon request}
% NOTE (CJ): keep names/emails out of this public repo. Maintain a separate
% one-page references list (name, title, relationship, contact) to send when
% a school asks, after confirming each person is willing.

\vfill
\begin{center}
	\footnotesize \textit{Updated \today}
\end{center}

\end{document}
